\section{Experiments}



In this section, we evaluate the proposed PRIME framework on both computer vision (CV) and natural language processing (NLP) tasks. All the experiments are executed on a single NVIDIA A100 GPU with PyTorch 2.10.0 and CUDA 12.8. 






\begin{table*}[t!]
  \centering
  \caption{Performance comparison of different methods on the ImageNet-1K dataset. ``Bits'' denotes the total number of bits transmitted across all timesteps (i.e., the cumulative bit-width of the spike over $T$ steps), reflecting the inter-layer data transmission cost during inference.}
  \label{tab:comparison}
  \resizebox{\textwidth}{!}{
  \begin{tabular}{llccccccc}
    \toprule
    Method & Arch. & Param.(M) & ANN Acc.(\%) & $T$ & Bits & SNN Acc.(\%) \\
    \midrule
    TS\cite{TS}
      & VGG-16         & 138.0    & 72.40 & 64           & 64          & 70.97~(-1.43) \\
    SNM\cite{SNM}
      & VGG-16         & 138.0    & 73.18 & 64           & 128         & 71.50~(-1.68) \\
    \cmidrule(lr){2-7}
    \multirow{2}{*}{QCFS\cite{QCFS}}
      & ResNet-34      & 21.8   & 74.32 & 64           & 64          & 72.35~(-1.97) \\
      & VGG-16         & 138.0    & 74.29 & 64           & 64          & 72.85~(-1.44) \\
    \cmidrule(lr){2-7}
    \multirow{2}{*}{SRP\cite{SRP}}
      & ResNet-34      & 21.8   & 74.32 & 64        & 64       & 68.61~(-5.71) \\
      & VGG-16         & 138.0    & 74.29 & 64        & 64       & 69.43~(-4.86) \\
    \cmidrule(lr){2-7}
    MST\cite{MST}
      & Swin-T         & 28.5   & 80.51 & 512     & 512    & 78.51~(-2.00) \\
    STA\cite{STA}
      & ViT-B/32       & 86.0     & 83.60 & 256      & 256     & 82.79~(-0.81) \\
    \cmidrule(lr){2-7}
    \multirow{3}{*}{SpikeZIP-TF\cite{spikeziptf}}
      & SViT-S-32 & 22.1  & 81.59 & 64           & 128         & 81.45~(-0.14) \\
      & SViT-B-32 & 86.6  & 82.83 & 64           & 128         & 82.71~(-0.12) \\
      & SViT-L-32 & 304.3 & 83.86 & 64           & 128         & 83.82~(-0.04) \\
    \cmidrule(lr){2-7}
    \multirow{2}{*}{ECMT\cite{ECMT}}
      & ViT-S/16       & 22.0     & 78.04 &10        & 40      & 77.07~(-0.97) \\
      & EVA-G          & 1074.0   & 89.62 &8         & 32      & 89.40~(-0.22) \\
    \cmidrule(lr){2-7}
    \makecell[l]{Adaptive fission\cite{adaptive}\\ + QCFS\cite{QCFS}}
      & ResNet-34      & 21.8   & 74.32 & 32           & 64          & 69.37~(-4.95) \\
    \cmidrule(lr){2-7}
    \multirow{4}{*}{DCGS\cite{dcgs}}
      & VGG-4/1        & 138.0    & 73.25 &  8         &  24      &  73.17~(-0.08) \\
      & ResNet34-4/1   & 21.8   & 76.42 & 8         &  24      &  76.04~(-0.38) \\
      & ViT-L-8/4      & 304.3  & 85.84 &  4         &  16       &  85.45~(-0.39) \\
      & EVA02-L-8/4    & 305.1  & 90.05 &  4         &  16       & 89.72~(-0.33) \\
    \midrule
    \multirow{10}{*}{\textbf{PRIME~(ours)}}
      & VGG-16         & 138.0    & 71.12 & \textbf{2} $(D{=}3)$  & \textbf{13}          & \textbf{71.08~(-0.04)} \\
      & ResNet-34      & 21.8   & 76.09 & \textbf{2} $(D{=}3)$  & \textbf{13}          & \textbf{76.13~(-0.00)} \\
      & ViT-S          & 22.1   & 78.04 & \textbf{2} $(D{=}3)$  & \textbf{13}          & \textbf{78.06~(-0.00)} \\
      & ViT-B          & 86.6   & 82.91 & \textbf{2} $(D{=}3)$  & \textbf{13}          & \textbf{82.87~(-0.04)} \\
      & ViT-L          & 304.3  & 84.88 & \textbf{2} $(D{=}3)$  & \textbf{13}          & \textbf{84.89~(-0.00)} \\
      & EVA02-T        & 5.8    & 80.62 & \textbf{2} $(D{=}3)$  & \textbf{13}          & \textbf{80.63~(-0.00)} \\
      & EVA02-S        & 22.1   & 85.73 & \textbf{2} $(D{=}3)$  & \textbf{13}          & \textbf{85.70~(-0.03)} \\
      & EVA02-B        & 87.1   & 88.26 & \textbf{2} $(D{=}3)$  & \textbf{13}          & \textbf{88.27~(-0.00)} \\
      & EVA02-L        & 305.1  & 89.63 & \textbf{2} $(D{=}3)$  & \textbf{13}          & \textbf{89.63~(-0.00)} \\
      & EVA02-G        & 1013.0   & 89.62 & \textbf{2} $(D{=}3)$  & \textbf{13}          & \textbf{89.58~(-0.04)} \\
    \bottomrule
  \end{tabular}
  }
\end{table*}

\subsection{Evaluation on Computer Vision Tasks}

For CV tasks, we validate the effectiveness of PRIME on the classic CNN and Transformer architectures, including VGG, ResNet, ViT, and EVA, evaluated on the ImageNet-1K~\cite{Imagenet} and CIFAR~\cite{cifar} (provided in Appendix~\ref{overall_cifars}, Figure~\ref{fig:cifar}) datasets. We follow the experimental settings of~\cite{ECMT}.










The experimental results on ImageNet-1K are presented in Table~\ref{tab:comparison}. \hwrevised{The results show that} PRIME consistently achieves lower conversion error across both CNN and Transformer architectures. Compared with the classic baseline DCGS~(\cite{dcgs}), PRIME uses only two timesteps to achieve lower accuracy degradation on VGG and ResNet ($0.04\%$ vs. $0.08\%$, $0.00\%$ vs. $0.38\%$). It demonstrates more pronounced advantages on ViT and EVA architectures ($0.00\%$ vs. $0.39\%$, $0.00\%$ vs. $0.33\%$). In summary, PRIME is the closest to lossless ANN-to-SNN conversion among the existing frameworks in Table \ref{tab:comparison}, while further reducing the inter-layer bit-width during hardware deployment ($13$ bits vs. $16$ bits). (\hwrevised{For additional analysis, the full  results across all $(T, D)$ configurations are provided in Appendix~\ref{overall_imagenet} (Table~\ref{tab:appendix_imagenet}). )}





































































\subsection{Evaluation on Natural Language Tasks}

For NLP tasks, we conduct the evaluations on Qwen, Llama, and OPT series models~\cite{qwen3,llama2,zhang2022opt}. The evaluations \hwrevised{perform} on the \hwrevised{primary NLP} benchmarks, including MMLU, PIQA, ARC-C, ARC-E, Winogrande, and Hellaswag~\cite{mmlu,piqa,wiqa, arc}. The evaluation process is implemented based on the work of~\cite{eval-harness}. We evaluate our method on the common sense reasoning tasks by converting LLAMA2-7b and LLAMA2-13b models with ultra-low latency ($T=2$). Following common practices in SNN conversion, we apply the same architectural configurations for our method. 
The results are presented in Table~\ref{tab:commonsense} and indicate that \hwrevised{ our method maintains an accuracy loss of less than 0.3$\%$ on the aforementioned NLP benchmarks.}
To the best of our knowledge, this is the first ANN-to-SNN conversion framework that consistently attains $<$0.3\% accuracy degradation on common-sense reasoning benchmarks, with only 2 timesteps. 
More results under different settings and datasets are shown in Appendix~\ref{overall_mmlu} (Figure~\ref{fig:mmlu}), Appendix~\ref{overall_llama} (Table~\ref{tab:appendix_llama}), and Appendix~\ref{overall_qwen3} (Table~\ref{tab:appendix_qwen3_mmlu}).
Notably, for other bit precision, our proposed method can also be applied effectively. 
The conversion results for FP8- and INT4-quantized LLMs are provided in Appendix~\ref{sec:fp8_int4}.








\begin{table*}[htbp]
  \centering
  \caption{Performance Comparison on Common Sense Reasoning Tasks (\%). ``Bits'' denotes the total bits transmitted across all timesteps, reflecting the data transmission cost during inference.}
  \label{tab:commonsense}
  \resizebox{\textwidth}{!}{
  \begin{tabular}{llcccccccccc}
    \toprule
    Method & Model & Spike & $T$ & Bits & PIQA & ARC-c\_n & ARC-e\_n & Winog. & HellS.\_n \\
    \midrule

    \multirow{2}{*}{LLAMA}
      & LLAMA2-7b  & $\times$ & N/A          & 16 & 78.45 & 45.73 & 74.87 & 69.38 & 76.28 \\
      & LLAMA2-13b & $\times$ & N/A          & 16 & 80.52 & 48.89 & 76.56 & 72.22 & 79.71 \\

    \multirow{2}{*}{SpikeLLM\cite{spikellm}}
      & LLAMA2-7b  & $\checkmark$ & 2        & 32 & 65.67 & 28.41 & 41.88 & 52.80 & 49.87 \\
      & LLAMA2-13b & $\checkmark$ & 2        & 32 & 67.68 & 30.38 & 44.53 & 57.22 & 58.02 \\

    FAS\cite{fas}
      & OPT-7B     & $\checkmark$ & 16       & 32 & 73.23 & ---   & 64.73 & 60.38 & ---   \\

    \multirow{3}{*}{LAS\cite{las}}
      & OPT-7B     & $\checkmark$ & 16       & 64 & 76.22 & ---   & 65.95 & 65.75 & ---   \\
      & OPT-13B    & $\checkmark$ & 16       & 64 & 76.28 & ---   & 67.38 & 65.51 & ---   \\
      & OPT-66B    & $\checkmark$ & 16       & 64 & 78.67 & ---   & 71.25 & 68.27 & ---   \\

    \cmidrule{1-10}

    \multirow{2}{*}{PRIME~(ours)}
      & LLAMA2-7b  & $\checkmark$ & \textbf{2 $(D{=}3)$} & \textbf{13} & \textbf{78.67} & \textbf{45.65} & \textbf{74.75} & \textbf{69.06} & \textbf{76.18} \\
      & LLAMA2-13b & $\checkmark$ & \textbf{2 $(D{=}3)$} & \textbf{13} & \textbf{80.58} & \textbf{49.15} & \textbf{76.47} & \textbf{72.14} & \textbf{79.69} \\

    \bottomrule
  \end{tabular}
  }
\end{table*}























\subsection{Ablation Study}

By default, we adopt the full FP16 exponent bit-width (5 bits) to set the threshold range, which preserves conversion accuracy. To evaluate the trade-off between efficiency and accuracy, we reduce the exponent bit-width to 4 and report the results in Table~\ref{tab:ablation2}. Across all $(T, D)$ configurations, the accuracy drop is marginal, while the number of active thresholds and the inter-layer bit-transmission cost is substantially reduced. This confirms that the proposed threshold-range optimization offers a practical knob for resource-constrained deployment with negligible performance sacrifice.
\hwrevised{More} ablation results on  the query spiking strategy in Transformer attention are provided in Appendix~\ref{sec:query_ablation}.


\begin{table*}[htbp]
  \centering
  \scriptsize  
  \caption{Ablation study on lightweight threshold-range optimization via exponent-bit 
  width reduction (\texttt{width=4}) on ImageNet-1K, evaluated across different 
  $(T, D)$ settings. \texttt{Default} refers to the full fp16 exponent-bit width 
  , with spike conversion applied to $k$ and $v$ only in Transformer 
  blocks.}
  \label{tab:ablation2}
  \resizebox{\textwidth}{!}{
  \begin{tabular}{cccccccccc}
    \toprule
    \multirow{2}{*}{Model} & \multirow{2}{*}{ANN} & \multirow{2}{*}{D} 
    & \multicolumn{3}{c}{Default} 
    & \multicolumn{3}{c}{Exponent-bit width $= 4$} \\
    \cmidrule(lr){4-6} \cmidrule(lr){7-9}
    & & & $T=1$ & $T=2$ & $T=3$ & $T=1$ & $T=2$ & $T=3$ \\
    \midrule
    \multirow{3}{*}{ViT-S (patch16\_224)} 
    & \multirow{3}{*}{78.04} 
    & 1 & 48.20 & 76.27 & 77.80 & 55.26 & 76.29 & 77.47 \\
    & & 2 & 73.78 & 77.94 & 78.03 & 73.94 & 77.66 & 77.65 \\
    & & 3 & 77.01 & 78.06 & 78.03 & 76.92 & 77.66 & 77.65 \\
    \midrule
    \multirow{3}{*}{ViT-B (patch16\_224)}
    & \multirow{3}{*}{82.91} 
    & 1 & 74.00 & 82.21 & 82.74 & 75.80 & 81.30 & 81.77 \\
    & & 2 & 80.89 & 82.77 & 82.89 & 80.30 & 81.74 & 81.88 \\
    & & 3 & 82.41 & 82.87 & 82.91 & 81.38 & 81.87 & 81.87 \\
    \midrule
    \multirow{3}{*}{ViT-L (patch16\_224)} 
    & \multirow{3}{*}{84.88} 
    & 1 & 80.89 & 84.46 & 84.79 & 81.14 & 83.81 & 84.31 \\
    & & 2 & 84.03 & 84.80 & 84.89 & 83.25 & 84.31 & 84.32 \\
    & & 3 & 84.45 & 84.89 & 84.88 & 84.01 & 84.30 & 84.30 \\
    \bottomrule
  \end{tabular}
  }
\end{table*}



















